
\needspace{4\baselineskip}
\color{uniblue}{\subsection{Предупредительная сигнализация устройства (ПС)\label{ps:punkt}}}
\color{black}

\begin{enumerate}[label=\arabic{section}.\arabic{subsection}.\arabic*, labelsep=4pt, leftmargin=0pt, itemindent=57pt]

\item	
Предупредительная сигнализация предназначена для информирования оперативного персонала при отклонениях от нормального режима работы оборудования энергообъекта.

\item
В устройстве предупредительная сигнализация представлена функцией «ПС».

\item
Функция принимает сигналы срабатывания защитных и сервисных функций устройства и формирует следующие сигналы:
\begin{itemize}
	\item импульсный сигнал предупредительной сигнализации «ПС: Пуск (имп)»;
	\item длительный сигнал предупредительной сигнализации «ПС: Пуск».
\end{itemize}

\item
Выходные сигналы функции могут быть использованы для дальнейшего конфигурирования (например, на выходные реле) и применяться в цепях участковых и индивидуальных табло, сигнальных ламп, сигнальных реле, звуковой сигнализации и пр.

\item
Введенный ключ «Вывод терминала» блокирует формирование выходных сигналов.

\item
Сигналы из логики СС: «Вых.цепи разобраны», «БИ выведены», «Неисправность КА», «Прев. вр.пер. КА», «Общ.внеш.сигнал» контролируются соответствующими уставками SGF1\ldots SGF7.

\item
Инвертированное состояние дискретных сигналов при введенных уставках SGF3 «Контроль ОТ цепей ОБР», SGF6 «Полож.пер.ав.дебл.» формирует работу логики ПС.

\item
Логика работы функции выполнена в соответствии с рисунком приложения \ref{app:ps}. В таблице \ref{ps:tbl1} приведены параметры, необходимые для настройки функционального блока «ПС». 

\small
\begin{longtable}{|>{\centering\arraybackslash}m{5.4cm}|>{\centering\arraybackslash}m{3.6cm}|>{\centering\arraybackslash}m{4.2cm}|>{\centering\arraybackslash}m{1.8cm}|>{\centering\arraybackslash}m{1cm}|}
\caption{Параметры для настройки функции «ПС»\hfill\vspace{-0.5\baselineskip}}\label{ps:tbl1}\\ 
\hline
\rowcolor{gray!30}
Параметр (Параметр на ИЧМ) & Условное обозначение на схеме & Значение/ Диапазон & Единица измерения & Шаг \\ 
\hline
\endfirsthead
%\caption{\emph{Продолжение\hfill\vspace{-0.5\baselineskip}}} \\ % Отступ обозначения таблицы от таблицы и выравнивание надписи слева
\caption*{\hspace{3pt}\emph{Продолжение таблицы \ref{ps:tbl1}\hfill\vspace{-0.5\baselineskip}}} \\ % сделано по ГОСТ 2.105 п.6.8.7
\hline
\rowcolor{gray!30}
Параметр (Параметр на ИЧМ) & Условное обозначение на схеме & Значение/ Диапазон & Единица измерения & Шаг \\ 
%\hline
\endhead
%\hline
\endfoot
\hline
\endlastfoot
%==+t1*CALH1|VT_LVALH
\centering Контроль выходных цепей (Контроль\_цепей) & \centering SGF1 & \centering 0 = Не предусмотрено\\1 = Предусмотрено & \centering -- & \centering \arraybackslash -- \\
\hline
\centering Контроль выведенного положения БИ (Контроль\_вывода\_БИ) & \centering SGF2 & \centering 0 = Не предусмотрено\\1 = Предусмотрено & \centering -- & \centering \arraybackslash -- \\
\hline
\centering Контроль сигнализации опер.тока (Контроль\_ОТ\_сигн) & \centering SGF3 & \centering 0 = Не предусмотрено\\1 = Предусмотрено & \centering -- & \centering \arraybackslash -- \\
\hline
\centering Контроль неисправности КА (Контроль\_неиспр\_КА) & \centering SGF4 & \centering 0 = Не предусмотрено\\1 = Предусмотрено & \centering -- & \centering \arraybackslash -- \\
\hline
\centering Контроль превышения времени переключения КА (Контр.прев.вр.пер.) & \centering SGF5 & \centering 0 = Не предусмотрено\\1 = Предусмотрено & \centering -- & \centering \arraybackslash -- \\
\hline
\centering Контроль положения переключателя аварийной деблокировки (Контр.пол.ав.дебл.) & \centering SGF6 & \centering 0 = Не предусмотрено\\1 = Предусмотрено & \centering -- & \centering \arraybackslash -- \\
\hline
\centering Контроль общего внешнего сигнала (Контр.общ.вн.сигн.) & \centering SGF7 & \centering 0 = Не предусмотрено\\1 = Предусмотрено & \centering -- & \centering \arraybackslash -- \\
%===t1
\end{longtable}
\normalsize

\end{enumerate}